\chapter{Introduction}\label{ch:intro}

\section{What GLoBES.jl does}

A long-baseline neutrino experiment has a source of neutrinos (an accelerator
beam, a neutrino factory, a beta beam or a nuclear reactor), a baseline of
$L$ kilometres through the Earth along which the neutrinos oscillate, and a
detector that counts neutrinos of certain flavours in bins of reconstructed
energy. GLoBES.jl computes
\begin{itemize}
  \item three-flavour oscillation probabilities in matter of any density profile
        (Chapter~\ref{ch:prob});
  \item the expected event rates in every energy bin, for signals and
        backgrounds, including the energy resolution of the detector
        (Chapter~\ref{ch:rates});
  \item the $\Delta\chi^2$ of a hypothesis compared with simulated data, with
        systematic uncertainties (Chapter~\ref{ch:chi2});
  \item the $\Delta\chi^2$ minimised over any set of oscillation parameters,
        with external information on them (``projections'' or
        ``marginalisation'', Chapter~\ref{ch:fits}), which is how GLoBES treats
        correlations and degeneracies.
\end{itemize}
Experiments are described in text files written in AEDL
(Part~\ref{part:aedl}). Their flux, cross section, energy resolution,
efficiency and systematics can be changed at runtime
(Chapter~\ref{ch:runtime}), and the oscillation probabilities can be replaced by
those of any non-standard model (Chapter~\ref{ch:nonstandard}).

\section{GLoBES.jl and GLoBES}

GLoBES.jl started as a line-by-line port of the C library GLoBES~3.0.11 and was
tested against it: the test suite compares probabilities, event rates and
$\chi^2$ values of thirteen experiment setups with values computed by the C
library. In the process, a number of bugs of the C library came to light. GLoBES.jl
fixes them, so that some results differ from C GLoBES; Chapter~\ref{ch:cdiffs}
lists every difference and why it exists.

GLoBES.jl has two interfaces:
\begin{description}
  \item[The high-level API] (Chapters~\ref{ch:tour}--\ref{ch:nonstandard} and
    \ref{ch:api}) uses named parameters (\code{OscParams(θ13=0.15, ...)}),
    experiment objects, keyword arguments, return values and exceptions. This
    is the interface to use for new programs.
  \item[The \code{glb*} functions] (Chapter~\ref{ch:compat}) have the names and
    arguments of the C functions, so that C GLoBES programs can be translated
    line by line. They work like the C library: errors are printed and the
    functions return error codes.
\end{description}
Both work on the same data and can be mixed.

\section{Installation}

GLoBES.jl needs Julia~1.10 or later. Its only dependency besides the standard
library is \code{SpecialFunctions}. Install it from its directory with
\begin{julia}
julia> using Pkg
julia> Pkg.develop(path="/path/to/GLoBES.jl")
\end{julia}
or start Julia with \code{julia --project=/path/to/GLoBES.jl}. Then
\begin{julia}
using GLoBES
\end{julia}
loads it. The test suite is run with \code{julia --project=. test/runtests.jl}
in the package directory.

\paragraph{Where AEDL files are found.} \code{Experiment("NOvA.glb")} looks for
the file in the current directory and then in the directories listed in the
environment variable \code{GLB\_PATH} (separated by \code{:}). Files that an AEDL
file includes, and its flux and cross section files, are looked up in the
current directory, next to the AEDL file, and in \code{GLB\_PATH}. The
experiments shipped with GLoBES.jl are in \file{test/data} and \file{examples};
Appendix~\ref{app:files} lists them.

\section{How this manual is organised}

Chapter~\ref{ch:tour} is a guided tour through a complete program. The
following chapters explain each step in detail together with the formulas
used. Sections marked with $^\ast$ describe advanced features that most users
will not need.
